\section*{Bloque B. Aplicaciones cuantitativas y optimización}

\noindent\textit{Preguntas 31 a 65. Requieren resolver el problema del monopolista, calcular precios, cantidades, elasticidades, índices de Lerner, pérdidas de bienestar y esquemas de discriminación y regulación.}

\noindent\textit{Considere un monopolio con demanda inversa $P=100-Q$ y costo total $C(Q)=20Q$, de modo que $CMg=20$. Las preguntas 31 a 45 se refieren a este caso, salvo indicación contraria.}

\q{La cantidad que maximiza el beneficio del monopolista es:}
\begin{opts}
    \item 20
    \item 30
    \item 40
    \item 50
\end{opts}

\q{El precio de monopolio es:}
\begin{opts}
    \item 40
    \item 60
    \item 80
    \item 100
\end{opts}

\q{El beneficio económico del monopolista es:}
\begin{opts}
    \item 800
    \item 1200
    \item 1600
    \item 2000
\end{opts}

\q{En el óptimo del monopolista, el ingreso marginal es:}
\begin{opts}
    \item 10
    \item 20
    \item 40
    \item 60
\end{opts}

\q{El índice de Lerner en el equilibrio del monopolista es aproximadamente:}
\begin{opts}
    \item 0.67
    \item 0.33
    \item 0.50
    \item 1.00
\end{opts}

\q{La elasticidad precio de la demanda, en valor absoluto, en el equilibrio del monopolista es:}
\begin{opts}
    \item 0.67
    \item 1.00
    \item 2.00
    \item 1.50
\end{opts}

\q{El excedente del consumidor en el equilibrio del monopolista es:}
\begin{opts}
    \item 400
    \item 800
    \item 1600
    \item 3200
\end{opts}

\q{La pérdida de bienestar del monopolio respecto de la solución competitiva es:}
\begin{opts}
    \item 400
    \item 600
    \item 1200
    \item 800
\end{opts}

\q{La cantidad socialmente óptima, donde $P=CMg$, es:}
\begin{opts}
    \item 40
    \item 80
    \item 100
    \item 60
\end{opts}

\q{El precio socialmente óptimo, donde $P=CMg$, es:}
\begin{opts}
    \item 40
    \item 60
    \item 20
    \item 0
\end{opts}

\q{El excedente total en la solución competitiva es:}
\begin{opts}
    \item 1600
    \item 2400
    \item 800
    \item 3200
\end{opts}

\q{El excedente total en el equilibrio del monopolio es:}
\begin{opts}
    \item 2400
    \item 3200
    \item 1600
    \item 800
\end{opts}

\q{Si el monopolista aplica discriminación de precios de primer grado, el excedente del consumidor es:}
\begin{opts}
    \item 800
    \item 0
    \item 1600
    \item 3200
\end{opts}

\q{Con discriminación de precios de primer grado, el beneficio del monopolista es:}
\begin{opts}
    \item 3200
    \item 2400
    \item 1600
    \item 800
\end{opts}

\q{Con discriminación de precios de primer grado, la cantidad producida es:}
\begin{opts}
    \item 40
    \item 80
    \item 100
    \item 60
\end{opts}

\q{El monopolista aplica discriminación de tercer grado entre dos mercados: $P_1=100-Q_1$ y $P_2=80-Q_2$, con $CMg=20$. La cantidad óptima en el mercado 1 es:}
\begin{opts}
    \item 40
    \item 30
    \item 20
    \item 50
\end{opts}

\q{En la discriminación de tercer grado anterior, la cantidad óptima en el mercado 2 es:}
\begin{opts}
    \item 40
    \item 30
    \item 20
    \item 50
\end{opts}

\q{En el mercado 2, el precio de equilibrio es:}
\begin{opts}
    \item 60
    \item 40
    \item 30
    \item 50
\end{opts}

\q{En discriminación de tercer grado, el grupo con demanda más elástica paga:}
\begin{opts}
    \item un precio menor.
    \item un precio mayor.
    \item el mismo precio que el otro grupo.
    \item un precio igual a cero.
\end{opts}

\q{En discriminación de tercer grado, el grupo con demanda menos elástica paga:}
\begin{opts}
    \item un precio menor.
    \item un precio mayor.
    \item el mismo precio que el otro grupo.
    \item un precio igual a cero.
\end{opts}

\q{Si la elasticidad de la demanda es 2 y el costo marginal es 10, el precio de monopolio según la regla de markup es:}
\begin{opts}
    \item 10
    \item 20
    \item 30
    \item 40
\end{opts}

\q{Si $|\varepsilon|=4$ y $CMg=30$, el precio de monopolio es:}
\begin{opts}
    \item 30
    \item 35
    \item 45
    \item 40
\end{opts}

\q{Si $P=50$ y $CMg=30$, el índice de Lerner es:}
\begin{opts}
    \item 0.2
    \item 0.3
    \item 0.4
    \item 0.6
\end{opts}

\q{Si el índice de Lerner es 0.5 y $P=80$, el costo marginal es:}
\begin{opts}
    \item 40
    \item 30
    \item 50
    \item 60
\end{opts}

\q{Considere un monopolio natural con $C(Q)=100+10Q$ y demanda $P=50-Q$. La cantidad que maximiza el beneficio es:}
\begin{opts}
    \item 10
    \item 15
    \item 20
    \item 25
\end{opts}

\q{En el monopolio natural anterior, el precio de monopolio es:}
\begin{opts}
    \item 20
    \item 40
    \item 10
    \item 30
\end{opts}

\q{En el monopolio natural anterior, el beneficio del monopolista es:}
\begin{opts}
    \item 100
    \item 200
    \item 400
    \item 300
\end{opts}

\q{Si el regulador fija $P=CMg=10$ en el monopolio natural anterior, la pérdida de la empresa es:}
\begin{opts}
    \item 0
    \item 100
    \item 200
    \item 300
\end{opts}

\q{Si el regulador fija un precio igual al costo medio, el beneficio económico de la empresa es:}
\begin{opts}
    \item positivo.
    \item negativo.
    \item máximo.
    \item cero.
\end{opts}

\q{En un monopolio natural no regulado, el precio de monopolio es:}
\begin{opts}
    \item igual al costo marginal.
    \item menor que el costo medio.
    \item mayor que el costo medio.
    \item igual a cero.
\end{opts}

\q{Si se aplica un precio igual al costo marginal en un monopolio natural, se requiere:}
\begin{opts}
    \item un subsidio para cubrir las pérdidas.
    \item un impuesto específico al consumo.
    \item un precio mínimo superior al CMe.
    \item ninguna intervención regulatoria adicional.
\end{opts}

\q{Una tarifa en dos partes con cargo fijo permite:}
\begin{opts}
    \item suprimir por completo los costos fijos operativos.
    \item establecer el cargo variable por unidad en cero.
    \item ampliar deliberadamente la pérdida de eficiencia social.
    \item cubrir costos fijos manteniendo precio cercano al CMg.
\end{opts}

\q{En el caso base, la diferencia entre el precio de monopolio y el precio competitivo es:}
\begin{opts}
    \item 20
    \item 30
    \item 40
    \item 60
\end{opts}

\q{En el caso base, el beneficio del monopolista también puede interpretarse como:}
\begin{opts}
    \item 1600
    \item 800
    \item 2400
    \item 3200
\end{opts}

\q{La pérdida de eficiencia del monopolio tiende a ser mayor cuando la demanda es:}
\begin{opts}
    \item perfectamente elástica.
    \item menos elástica.
    \item horizontal.
    \item de elasticidad infinita.
\end{opts}
